\subsection{赋范对合代数}

定义 2. --- 赋范对合代数是装备了一个对合 \(x \mapsto x^{*}\) 的赋范代数 A，满足 \(\|x^{*}\| = \|x\|\) 对一切 \(x\) 成立。若 A 是 Banach 代数，则称 A 为对合 Banach 代数。

例. --- 1) 设 X 是局部紧拓扑空间。X 上连续有界复函数组成的 Banach 代数 \(\mathcal{C}_{b}(\mathrm{X})\)，装备范数 \(\|f\| = \sup_{x\in\mathrm{X}}|f(x)|\) 与对合 \(f\mapsto\overline{f}\)，是一个对合 Banach 代数。在无穷远处趋于 0 的连续函数组成的子代数 \(\mathcal{C}_{0}(\mathrm{X})\) 是 \(\mathcal{C}_{b}(\mathrm{X})\) 的一个闭对合子代数。

\begin{enumerate}
\def\labelenumi{\arabic{enumi})}
\setcounter{enumi}{1}
\item
  设 X 是局部紧拓扑空间，\(\mu\) 是 X 上的一个正测度。对合代数 \(\mathrm{L}^{\infty}(\mathrm{X},\mu)\) (I, p.~99, 例 2) 是对合 Banach 代数，因为 \(|f| = |\overline{f}|\) 对每个元素 \(f \in \mathrm{L}^{\infty}(\mathrm{X},\mu)\) 成立。
\item
  复 Hilbert 空间 E 的连续自同态组成的对合代数 \(\mathcal{L}(\mathrm{E})\) (I, p.~99, 例 3)，装备范数
\end{enumerate}

\[
\| u\| = \sup_{\substack{x\in \operatorname{E}\\ \| x\| \leqslant 1}}\| u(x)\|
\]

（EVT, III, p.~14）后，是一个对合 Banach 代数 (EVT, V, p.~37, 命题 1)。

\begin{enumerate}
\def\labelenumi{\arabic{enumi})}
\setcounter{enumi}{3}
\item
  局部紧群上有界测度组成的对合代数 \(\mathcal{M}^{1}(G)\) (I, p.~99, 例 4)，装备通常的范数 (I, p.~19, 例 6)，是对合 Banach 代数。设 \(\nu\) 是 G 上的一个左 Haar 测度。对合 Banach 代数 \(\mathrm{L}^{1}(\mathrm{G},\nu)\) 恒等同于 \(\mathcal{M}^{1}(\mathrm{G})\) 的一个闭子代数。
\item
  设 \((\mathrm{A}_{i})\) 是一族对合赋范代数。设 A 是 \(A_{i}\) 的乘积赋范代数 (I, p.~15, 第 1 小节)。代数 A 装备对合 \((x_{i})^{*} = (x_{i}^{*})\) 后是一个对合赋范代数。若每个代数 \(A_{i}\) 都是对合 Banach 代数，则 A 是对合 Banach 代数。我们称 A 为 \(A_{i}\) 的乘积对合赋范代数（相应地，乘积对合 Banach 代数）。
\item
  设 A 是一个赋范对合代数，\(\widetilde{\mathrm{A}}\) 是由 A 添加单位元而得的赋范代数。代数 \(\widetilde{\mathrm{A}}\) 装备第 2 小节中定义的对合后是一个赋范对合代数。若 A 是对合 Banach 代数，则 \(\widetilde{\mathrm{A}}\) 也是对合 Banach 代数。
\end{enumerate}

若 A 是赋范对合代数，对合子代数的闭包是对合子代数。若 M ⊂ A，包含 M 的最小闭对合子代数称为由 M 生成的闭对合子代数；它是由 M ∪ M* 生成的子代数的闭包。若 M 由单个正规元组成，则由 M 生成的闭对合子代数是交换的，且其所有元素都是正规的。

类似地，若 A 是幺对合赋范代数且 M 是 A 的一个子集，包含 M 的最小闭幺对合子代数称为由 M 生成的闭幺对合子代数；它是由 M ∪ M* 生成的幺子代数的闭包。若 M 由单个正规元组成，则由 M 生成的闭幺对合子代数是交换的，且其所有元素都是正规的。

对合赋范代数对闭自伴双边理想的商、完备化以及对合赋范代数的反代数，都自然是赋范对合代数。

若 A 是赋范对合代数，A 的 Hermite 元组成的集合 \(A_{h}\) 是一个实赋范向量空间。

引理 4. --- 设 A 是一个赋范对合代数。对 A 上每个连续线性型 f，我们有 \(\|f^{*}\|=\|f\|\)。此外若 f 是 Hermite 的，则 \(\|f\|=\|f|A_{h}\|\)。

第一个断言由定义得出。对第二个断言，设 \(g\) 是 \(f\) 在 \(\mathrm{A}_h\) 上的限制。我们有 \(\| f \| \geqslant \| g \|\)。我们来证明反向不等式。对每个 \(\varepsilon > 0\)，存在 \(x \in \mathrm{A}\) 使 \(\| x \| \leqslant 1\) 且 \(|f(x)| \geqslant \|f\| - \varepsilon\)。用模为 1 的复数乘 x，可设 \(f(x) \geqslant 0\)。于是元素 \(\frac{1}{2}(x + x^{*})\) 属于 \(A_{h}\) 且范数 \(\leqslant 1\)。我们有

\[
\left| g \left(\frac {1}{2} (x + x ^ {*})\right) \right| = \frac {1}{2} | f (x) + f (x ^ {*}) | = f (x) \geqslant \| f \| - \varepsilon
\]

于是 \(\|g\|\geqslant\|f\|-\varepsilon\)。由此我们推出 \(\|g\|\geqslant\|f\|\)。

在下文中，我们将把 A 上连续的 Hermite 线性型与 \(A_{h}\) 上连续的实线性型恒同。

引理 5. --- 设 A 是一个对合 Banach 代数。

\begin{enumerate}
\def\labelenumi{\alph{enumi})}
\item
  对每个 \(x \in A\)，有 \(\exp(x)^{*} = \exp(x^{*})\)；
\item
  设 \(x \in \mathrm{A}_h\) 是一个 Hermite 元。那么 \(\exp(ix)\) 是酉元。
\end{enumerate}

事实上，由于 A 上的对合是连续的，我们有

\[
\exp (x) ^ {*} = \left(\sum_ {n = 0} ^ {\infty} \frac {x ^ {n}}{n !}\right) ^ {*} = \sum_ {n = 0} ^ {\infty} \frac {(x ^ {*}) ^ {n}}{n !} = \exp (x ^ {*})
\]

对一切 \(x \in A\) 成立 (I, p.~78, 公式 (9))。若 \(x \in A_{h}\)，则由此得

\[
\exp (i x) ^ {*} = \exp ((i x) ^ {*}) = \exp (- i x) = \exp (i x) ^ {- 1}
\]

(I, p.~78, 公式 (11))。

